\chapter{Cuatro ideas antes de seguir}\label{cap:conceptos}

Para usar la librería hacen falta cuatro ideas. Ninguna es difícil, pero conviene tenerlas claras: casi todos los errores de migración vienen de confundir alguna.

\section{Las entradas y salidas viajan en estructuras}
En un PLC Siemens, los sensores y actuadores llegan como bits en la \emph{imagen de proceso}: la entrada \cod{\%I1.0} es el primer bit del byte 1. Para no manejar bits sueltos, la librería agrupa los bits de cada módulo en un tipo de datos (UDT) con nombres: por ejemplo \cod{UDT\_Entradas\_Basico} tiene los miembros \cod{Boton\_marcha}, \cod{Boton\_paro}, etc. Después, \emph{una variable PLC} (un \emph{tag}) de ese tipo, situada en la dirección del módulo, hace que cada bit tenga su nombre.

\begin{idea}
El tag \cod{Entradas\_Basico} de tipo \cod{UDT\_Entradas\_Basico} en la dirección \cod{\%I1.0} \textbf{es} el byte de entradas 1 del PLC, visto con nombres. No hay que copiar nada: el nombre y el bit son la misma cosa.
\end{idea}

La librería usa cinco de estos tags comunes (los tienes en el apéndice~\ref{ap:direcciones}). Tu aplicación puede tener los suyos (\cod{UDT\_Entradas\_APP} y \cod{UDT\_Salidas\_APP}), que la librería no ve nunca.

\section{El palet lleva una memoria: el tag RFID}
Cada palet lleva un transpondedor (\emph{tag}) con 108 bytes de memoria. Es lo único que \textbf{viaja con el palet}: lo que una estación escribe en el tag, la siguiente lo lee. Por eso el tag es el idioma común de la línea.

El lector de cada estación (un SIMATIC RF210R conectado por IO-Link) no deja leer byte a byte. Funciona así:
\begin{itemize}
\item Se le da una \textbf{dirección de inicio} (un \emph{offset} en bytes dentro del tag) y una orden de leer o de escribir.
\item Lee o escribe siempre \textbf{un bloque de 28 bytes} a partir de esa dirección.
\item Termina cuando la señal \emph{proceso listo} baja y vuelve a subir, y la dirección que devuelve coincide con la que se pidió.
\end{itemize}

\begin{ojo}
Una escritura \textbf{siempre escribe los 28 bytes}, no solo los que tú quieres cambiar. Si escribes sin haber leído antes, los bytes que no tocas se quedan a cero y \emph{borras} lo que habían escrito las otras estaciones. Esta es la razón de ser de las máscaras (idea 3).
\end{ojo}

Hay un detalle más: el grupo acordó usar un único bloque de 28 bytes al principio del tag, a partir de la dirección que asigne el profesor (0 si el bloque es el primero). El reparto de esos 28 bytes está en el capítulo~\ref{cap:mapa}.

\section{La imagen del palet y las máscaras}
Para que tu aplicación no tenga que preocuparse del lector, la librería le entrega los datos del palet en una variable llamada \cod{Imagen\_palet}, un array de 28 bytes. Funciona como una fotocopia del bloque del tag con permisos:

\begin{figure}[ht]
\centering
\begin{tikzpicture}[x=1cm,y=1cm]
\node[hwcaja,minimum width=2.6cm,minimum height=1.5cm,font=\sffamily\small] (tag) at (0,0) {\textbf{Tag}\\[-1pt]\scriptsize 28 bytes del bloque};
\node[libcaja,minimum width=3.2cm,minimum height=1.5cm,font=\sffamily\small] (rf) at (4.4,0) {\cod{Est\_FB\_RFID}\\[-1pt]\scriptsize lee y escribe\\[-2pt]\scriptsize con las máscaras};
\node[caja,fill=white,minimum width=2.8cm,minimum height=1.5cm,font=\sffamily\small] (img) at (8.5,0) {\cod{Imagen\_palet}\\[-1pt]\scriptsize 28 bytes};
\node[appcaja,minimum width=2.6cm,minimum height=1.5cm,font=\sffamily\small] (app) at (12.3,0) {\textbf{Tu aplicación}};
\draw[flecha,<->] (tag) -- (rf);
\draw[flecha,<->] (rf) -- (img);
\draw[flecha,<->] (img) -- (app);
\node[etiqueta,above=2pt of $(tag.east)!0.5!(rf.west)$] {orden leer / escribir};
\node[etiqueta,below=14pt of rf,align=center] {\cod{Mascara\_lectura}: qué bytes llegan a la imagen\\\cod{Mascara\_escritura}: qué bytes vuelven al tag};
\end{tikzpicture}
\caption{El camino de los datos del palet.}
\end{figure}

\begin{itemize}
\item \textbf{Al leer}, a la imagen solo llegan los bytes que marca \cod{Mascara\_lectura}; el resto de la imagen queda a cero. Tu aplicación solo ve lo que necesita.
\item \textbf{Al escribir}, al tag solo van los bytes de la imagen que marca \cod{Mascara\_escritura}; los demás se devuelven al tag \emph{tal como se leyeron}. Tu aplicación no puede estropear los bytes de otras estaciones.
\end{itemize}

Cada máscara es un número de 32 bits en el que el bit $n$ corresponde al byte $n$ del bloque. En el capítulo~\ref{cap:mapa} se aprende a calcularlas.

\section{Tipos de la librería y copias maestras}\label{sec:tipos-copias}
Una librería global de TIA Portal guarda dos clases de cosas, que se usan de manera distinta:

\begin{description}[style=nextline,leftmargin=1.2cm]
\item[Tipos (\emph{Types}).] Tienen \textbf{versión}. Al arrastrarlos a tu proyecto quedan \emph{ligados} a la librería: no se editan en tu proyecto, y si sale una versión nueva se actualizan con un par de clics. Son las piezas que no debes tocar: los cuatro FB y los UDT comunes.
\item[Copias maestras (\emph{Master copies}).] Son copias normales, sin versión ni vínculo. Se arrastran a tu proyecto y son \textbf{tuyas}: las cambias como quieras. Son los puntos de partida: la plantilla de aplicación, el \cod{Main} de ejemplo, el bloque de datos y la tabla de variables.
\end{description}

\begin{idea}
Los tipos son \textbf{para usarlos}; las copias maestras son \textbf{para modificarlas}. Si algo tuyo tiene que cambiar, está en las copias maestras. Si algo de la librería tiene que cambiar, se pide un cambio en el tipo para todos.
\end{idea}

\section{Glosario}
\begin{center}
\small
\begin{tabularx}{\textwidth}{>{\bfseries}p{3.4cm}X}
\toprule
Palabra & Qué significa aquí \\
\midrule
Palet & La plataforma que lleva la pieza por la línea. Lleva un tag RFID. \\
Tag RFID & El transpondedor del palet; 108 bytes de memoria. La librería usa un bloque de 28. \\
Bloque & Los 28 bytes del tag que acordó el grupo (capítulo~\ref{cap:mapa}). \\
Imagen del palet & La copia en el PLC del bloque, con permisos (\cod{Imagen\_palet}). \\
Máscara & Número de 32 bits que dice qué bytes se leen o se escriben. \\
Contrato & Las ocho señales y la imagen que cumple tu aplicación para trabajar con la librería. \\
Aplicación & El bloque de función que es propio de cada estación. \\
FB / OB / DB & Bloque de función / bloque de organización (\cod{Main}) / bloque de datos. \\
UDT & Tipo de datos definido por el usuario (una estructura con nombres). \\
Tag (PLC) & Una variable del PLC con nombre y dirección. No confundir con el tag RFID. \\
Instancia & Cada llamada a un FB tiene su propia memoria, su bloque de datos de instancia. \\
Librería global & Archivo de TIA con tipos y copias maestras que se comparte entre proyectos. \\
\bottomrule
\end{tabularx}
\end{center}
